% mnras_template.tex 
\documentclass[letters,fleqn,usenatbib]{mnras}

% Import showyourwork magic
\usepackage{showyourwork}

% MNRAS is set in Times font. If you don't have this installed (most LaTeX
% installations will be fine) or prefer the old Computer Modern fonts, comment
% out the following line
%\usepackage{newtxtext,newtxmath}
% Depending on your LaTeX fonts installation, you might get better results with one of these:
\usepackage{mathptmx}
%\usepackage{txfonts}

% Use vector fonts, so it zooms properly in on-screen viewing software
% Don't change these lines unless you know what you are doing
\usepackage[T1]{fontenc}

% Allow "Thomas van Noord" and "Simon de Laguarde" and alike to be sorted by "N" and "L" etc. in the bibliography.
% Write the name in the bibliography as "\VAN{Noord}{Van}{van} Noord, Thomas"
\DeclareRobustCommand{\VAN}[3]{#2}
\let\VANthebibliography\thebibliography
\def\thebibliography{\DeclareRobustCommand{\VAN}[3]{##3}\VANthebibliography}


%%%%% AUTHORS - PLACE YOUR OWN PACKAGES HERE %%%%%

% Only include extra packages if you really need them. Avoid using amssymb if newtxmath is enabled, as these packages can cause conflicts. newtxmatch covers the same math symbols while producing a consistent Times New Roman font. Common packages are:
\usepackage{graphicx}	% Including figure files
\usepackage{amsmath}	% Advanced maths commands
\usepackage{siunitx}

%%%%%%%%%%%%%%%%%%%%%%%%%%%%%%%%%%%%%%%%%%%%%%%%%%

%%%%% AUTHORS - PLACE YOUR OWN COMMANDS HERE %%%%%

% Please keep new commands to a minimum, and use \newcommand not \def to avoid
% overwriting existing commands. Example:
%\newcommand{\pcm}{\,cm$^{-2}$}	% per cm-squared

%%%%%%%%%%%%%%%%%%%%%%%%%%%%%%%%%%%%%%%%%%%%%%%%%%

%%%%%%%%%%%%%%%%%%% TITLE PAGE %%%%%%%%%%%%%%%%%%%

% Title of the paper, and the short title which is used in the headers.
% Keep the title short and informative.
\title[The dispersal of compact protoplanetary discs]{The dispersal of compact protoplanetary discs}

% The list of authors, and the short list which is used in the headers.
% If you need two or more lines of authors, add an extra line using \newauthor
\author[G. Picogna et al.]{
Giovanni Picogna$^{1}$\thanks{E-mail: picogna@usm.lmu.de}
and Barbara Ercolano$^{1,2,3}$
\\
% List of institutions
$^{1}$Universit\"{a}ts-Sternwarte, Ludwig-Maximilians-Universit\"{a}t M\"{u}nchen, Scheinerstr. 1, D-81679 M\"{u}nchen, Germany\\
$^{2}$Excellence Cluster Origins, Boltzmannstrasse 2, D-85748 Garching bei M\"{u}nchen, Germany\\
$^{3}$Max-Planck-Institut fu\"{u}r Extraterrestrische Physik, Giessenbachstraße 1, 85748 Garching, Germany
}

% These dates will be filled out by the publisher
\date{Accepted XXX. Received YYY; in original form ZZZ}

% Enter the current year, for the copyright statements etc.
\pubyear{2025}

% Don't change these lines
\begin{document}
\label{firstpage}
\pagerange{\pageref{firstpage}--\pageref{lastpage}}
\maketitle

% Abstract of the paper
\begin{abstract}
Compact protoplanetary discs are becoming increasingly prominent in observations. Their dispersal pathways may differ substantially from those of extended discs. We aim to quantify the role of the disc outer radius in internal photoevaporation, provide a simple scaling relation for compact discs, and test whether the resulting evolutionary tracks reproduce the observed inside-out clearing of young stellar populations. We performed radiation-hydrodynamic simulations of X-ray-driven photoevaporation for discs with different outer radii, and derived the dependence of the total mass-loss rate on the critical radius. We incorporated this scaling into disc population synthesis models. We find that the surface mass-loss profiles are nearly independent of disc size, but their integrated wind rates are reduced according to the cumulative mass-loss rate distribution. Compact discs disperse more slowly, evolving via inside-out clearing consistent with observational diagnostics. Population synthesis shows that the scaled prescription reproduces both disc fractions and the persistence of highly accreting systems at late times ($t > 4$ Myr) better than unscaled models. Disc dispersal prescriptions must include the dependence on disc outer radius to capture the evolution of compact discs. The proposed scaling provides a simple, physically motivated correction that better predicts the growing observational evidence for compact discs and inside-out dispersal.
\end{abstract}

% Select between one and six entries from the list of approved keywords.
% Don't make up new ones.
\begin{keywords}
keyword1 -- keyword2 -- keyword3
\end{keywords}

%%%%%%%%%%%%%%%%%%%%%%%%%%%%%%%%%%%%%%%%%%%%%%%%%%

%%%%%%%%%%%%%%%%% BODY OF PAPER %%%%%%%%%%%%%%%%%%

\section{Introduction}\label{sec:introduction}

 Protoplanetary discs are the birthplaces of planets, and their lifetimes and dispersal pathways set essential constraints on planet formation models. Most discs dissipate within a few Myr, with only a minority surviving to \SI{\sim 10}{Myr} \citep{2001ApJ...553L.153H, 2010A&A...510A..72F}. Observations also show that dispersal proceeds preferentially from the inside-out: colour–colour evolution studies of young stellar populations \citep{2011MNRAS.410..671E, 2013MNRAS.428.3327K, 2015MNRAS.452.3689E} reveal clear signatures of inner clearing, inconsistent with outside-in erosion. Any successful theoretical model must therefore reproduce this inside-out behaviour.

 Mounting evidence points to magnetically driven winds and photoevaporation (EUV, X-ray, FUV) as the dominant agents of angular momentum removal and mass loss, respectively \citep{2017RSOS....470114E, 2023ASPC..534..465L, 2023ASPC..534..567P}. X-ray photoevaporation models in particular have been refined in recent years to yield dispersal prescriptions consistent with several independent observational diagnostics \citep{2016MNRAS.460.3472E, 2020MNRAS.496..223W, 2021A&A...648A.121F, 2023ApJ...955L..11R, 2023MNRAS.526L.105E}.

 Most existing models, however, assume discs with outer radii of order hundreds of au - well beyond the gravitational radius $R_g$, defined as the location where the gas sound speed becomes comparable with the Keplerian speed - around 5 au for a Solar mass star, creating a natural "sweet spot" for photoevaporation \citep{1994ApJ...428..654H, 2003PASA...20..337L}. Yet compact discs, with outer radii comparable to or smaller than this scale, are increasingly recognised as common. ALMA surveys show that a substantial fraction of discs are indeed very compact \citep[e.g.][]{Trapman_2025}, with critical radii much smaller than typically assumed. Such compact configurations may result from internal evolution (mass and angular momentum loss via winds) or from external influences such as stellar fly-bys or external photoevaporation.

 In this Letter, we investigate the dispersal of compact discs by explicitly accounting for their finite outer radii in photoevaporation models. Using radiation-hydrodynamic simulations, we derive the dependence of the total wind mass-loss rate on the critical radius and implement this scaling into one-dimensional viscous evolution and population synthesis models. Our aim is to provide a simple prescription valid in the compact-disc regime and to test whether the resulting evolutionary pathways are consistent with observed disc fractions, accretion rates, and the observationally inferred inside-out dispersal of young stellar populations.

\section{Methods}\label{sec:methods}

We performed a series of radiation hydrodynamic simulations following the approach described in \citet{2019MNRAS.487..691P}, which we briefly summarise here for completeness.

\subsection{Radiative Transfer}

We focused on a \SI{1}{M_\odot} star, where the stellar properties were taken from \citet{2000A&A...358..593S}, assuming an age of 1 Myr and metallicity $Z=0.02$ without convective overshooting.
The initial disc structure was based on the gas densities and dust temperatures from the hydrostatic disc models of the D’Alessio Irradiated Accretion disc (\textsc{diad}) radiative transfer framework \citep{1998ApJ...500..411D, 1999ApJ...527..893D, 2001ApJ...553..321D, 2005ApJ...621..461D, 2006ApJ...638..314D}. These models provide the best fits to the median spectral energy distribution (SED) observed in Taurus.
To compute the thermal structure of the upper layers of the disc, we used the gas photoionisation and dust radiative transfer code \textsc{mocassin} \citep{2003MNRAS.340.1136E,2005MNRAS.362.1038E,2008ApJS..175..534E}.
This code self-consistently solves for heating and cooling processes under thermal equilibrium, allowing us to derive temperature prescriptions up to the maximum penetration depth of X-rays ($\sim 10^{22}$ pp/cm$^2$) for discs irradiated by X-ray and EUV stellar spectra.
The resulting temperature prescription depends on the spectral hardness as detailed in \citep{2021MNRAS.508.1675E}.
It relates the local gas temperature to both the column density towards the central star (from \SI{5e20}{pp.cm^{-2}} to \SI{2e22}{pp.cm^{-2}}), and the local ionisation parameter \citep{1969ApJ...156..943T} $\xi = L_X/(n r^2)$, where $L_X$ is the stellar X-ray luminosity, $n$ the local number density, and $r$ the spherical radius.

\subsection{Hydrodynamical model}

Having the temperature prescription at thermal equilibrium, we then performed hydrodynamical simulations using a modified version of the \textsc{pluto} code \citep{2007ApJS..170..228M}, as described in \citet{2019MNRAS.487..691P}.
In this setup, the \textsc{mocassin}-derived prescription was applied at column densities below the X-ray penetration depth (\SI{\sim 2e22}{pp.cm^{-2}}), while at larger column densities we assumed perfect thermal coupling between gas and dust, adopting the \textsc{diad} models.
Each simulation was evolved until the disc structure and wind streamlines reached steady state.
We explored $4$ different initial critical radii for the circumstellar disc: $10$, $50$, $100$, $200$ au, where we exponentially cut off the gas surface densities from the unperturbed disc.
We adopted a spherical coordinate system centred on the star.
The computational grid is logarithmically spaced in the radial direction with $500$ cells, ensuring enhanced resolution in the innermost disc regions where photoevaporation is most effective, while still extending to large radii ($R_\mathrm{out} = 300\,\mathrm{au}$) without excessive computational cost. In the polar direction, the grid is linearly spaced with $250$ cells between \SI{0.01}{rad} and $\pi/2$.
The system was evolved for $500-1000$ orbits at \SI{10}{au}.
A quasi-steady state is typically reached after a few hundred orbital periods, at which point the cumulative mass-loss rate and the gas streamlines in the wind stabilise.

\subsection{Disc population synthesis}

In order to model the long-term evolution of compact protoplanetary discs, we adopted a 1D viscous evolution code \citep{2015MNRAS.450.3008E} where the 1D surface density is evolved according to
\begin{equation}\label{eq:visc}
\frac{\partial\Sigma}{\partial t} = \frac{1}{R} \frac{\partial}{\partial R} \left[ 3R^{1/2} \frac{\partial}{\partial R} \left( \nu\Sigma R^{1/2} \right)\right] - \dot{\Sigma}_{\mathrm{w}}(R, t)\,,
\end{equation}
where the first term describes the disc viscous evolution \citep{1974MNRAS.168..603L} and the last term the mass loss due to photoevaporation derived from our hydrodynamical models \citep{2021MNRAS.508.3611P}
\begin{eqnarray}
  \label{eq:surf1}
  \dot{\Sigma}_w(R) &= \ln{(10)} \bigg(\frac{6\, a\, \ln{(R)}^5}{R\, \ln{(10)}^6} +
  \frac{5\, b\, \ln{(R)}^4}{R\, \ln{(10)}^5} +
  \frac{4\, c\, \ln{(R)}^3}{R\, \ln{(10)}^4} + \\ \nonumber
  &\frac{3\, d\, \ln{(R)}^2}{R\, \ln{(10)}^3} +
  \frac{2\, e\, \ln{(R)}}{R\, \ln{(10)}^2} + \\ \nonumber
  &\frac{f}{R\, \ln{(10)}}\bigg)
  \frac{\dot{\mathrm{M}}_w(R)}{2\pi\, R} \ [\SI{}{M_\odot \textrm{au}^{-2} yr^{-1}}]\,,
\end{eqnarray}
where 
\begin{equation}
\label{eq:surf2}
  \dot{\mathrm{M}}_w(R) = 10^{a\log{R}^6 + b\log{R}^5 + c\log{R}^4 + d\log{R}^3 + e\log{R}^2 + f\log{R} + g}
\end{equation}
is the cumulative mass-loss rate, and the parameters $a$, $b$, $c$, $d$, $e$, $f$, $g$ are stellar-mass dependant and given in \citet[][Table 2]{2021MNRAS.508.3611P}.
The discs are initialized with a self-similar solution (for $\gamma = 1$):
\begin{equation}
    \Sigma = \frac{M_0}{2 \pi r_1 R} T^{-3/2} \exp{\left[-\frac{(R/r_1)}{T}\right]}\,,
\end{equation}
where $M_0$ is the initial disc mass, $T=t/t_s$ is the dimensionless time with $t_s$ the viscous timescale, and $r_1$ is the disc characteristic or critical radius.

We performed a parameter space analysis by probing a large set of disc and stellar properties.
We assumed a stellar initial mass function following \citet{2001MNRAS.322..231K}, $\xi(m)\propto m^{-\alpha}$, where $\alpha = 1.3 \pm 0.5$ for $0.08 \leq m/M_\odot < 0.5$ and $\alpha = 2.3 \pm 0.3$ for $0.5 < m/M_\odot < 1$.
For the X-ray luminosities we followed the same approach outlined in \citet{2023MNRAS.526L.105E}, deriving an X-ray luminosity function in three stellar mass bins from a subsample of the Chandra Orion Ultra-deep Project \citep[COUP, cf.][]{2005ApJS..160..319G} as shown in \citet{2019ApJ...883..117K}.
For the disc properties, we derive the disc dust mass from observational constraint in young star forming regions \citep{2016ApJ...831..125P, 2017AJ....153..240A}:
\begin{equation}
    \log(M_\mathrm{dust} [M_\oplus]) = (1.2\pm0.2) + (1.8\pm0.4) \log(M_\star [M_\odot])
\end{equation}
with an associated uncertainty of $\delta = 0.9\pm0.1$.
We then consider a standard 100 conversion factor between the dust and the gas mass, and derive the critical radius by inverting the observed dependence of the critical radius as a function of the gas mass \citep{Trapman_2025}, $M_\mathrm{gas} \propto r_1^{1.7}$. Finally, the disc viscosity is the less constrained parameter, thus we linearly sampled its values between $\alpha = 10^{-4}$ and $10^{-2 }$. The resulting parameter space sampled is shown in Fig.~\ref{fig:pop_synth_corner_plots}.
\begin{figure*}
    \centering
    \includegraphics[width=17cm]{figures/Fig1.png}
    \caption{Parameter space sampled in the population synthesis. In panel a, stellar mass as a function of X-ray luminosity sampled. In panel b, critical radius as a function of disc mass sampled. The observationally derived fits from \citet{2007A&A...468..353G} and \citet{Trapman_2025} are overplotted for comparison.}
    \script{fig1.py}
    \label{fig:pop_synth_corner_plots}
\end{figure*}

\section{Results}\label{sec:results}

\subsection{Hydrodynamical models}

We computed the cumulative mass-loss rate from our hydrodynamical models, as shown in Fig.~\ref{fig:mass_loss_rate_cumulative}. The four cases exhibit very similar radial behaviours, with $\dot{M}_w(R)$ rising steeply in the inner disc and flattening towards large radii. This outcome is expected, since the incident radiation field at smaller radii is not directly affected by the details of the density distribution further out.
The main effect of reducing the critical radius is to truncate the radial extent of the mass-loss distribution (equation~\ref{eq:surf2}), thereby lowering the total integrated mass-loss rate.
In Fig.~\ref{fig:mass_loss_rate_cumulative} we show also that the complex fit of equation~\ref{eq:surf2} can be actually broken in two linear regimes, in order to find a simpler (and physically motivated) way to describe the wind mass-loss, where for radii comparable with the gravitational radius it grows as $\dot{M}_w \propto R^3$, while for larger radii like $\dot{M}_w \propto R^{0.5}$. This shows how important it is to consider the critical radius truncated mass-loss rate for compact discs, where the bulk of the disc fits inside the fast growth regime of the wind mass-loss rate.

We can now provide a prescription for the long-term evolution of compact discs. The normalization factor to be used is given by subtituting in equation~\ref{eq:surf2} the critical radius $r_1$
\begin{equation}\label{eq:mdotnorm}
    \dot{M}_{w, \mathrm{norm}} = \dot{M}_{w}(r_1)\,.
\end{equation}
The total mass-loss rate for a compact disc will then be:
\begin{equation}\label{eq:mdot}
    \dot{M}_{w} = \dot{M}_{w, \mathrm{norm}} \dot{M}_w(M_\star) \dot{M}_w(L_X)\,,
\end{equation}
where the mass-loss rate as a function of stellar mass (for a mean X-ray luminosity, \citet{2021MNRAS.508.3611P}) is
\begin{equation}\label{eq:mdotmstar}
    \dot{M}_w(M_\star) = 3.93\times10^{-8} \left(\frac{M_\star}{M_\odot}\right)
\end{equation}
and the dependence on the X-ray luminosity \citep{2021MNRAS.508.1675E} is 
\begin{equation}\label{eq:mdotlx}
    \begin{split}
    \dot{M}_w(L_X) &= \frac{\dot{M}_w(L_{X,\mathrm{soft}})}{\dot{M}_w(L_{X,\mathrm{soft, mean}})} \\
    &= 10^{a_L \left[\exp{\left(\frac{(\ln(\log(L_{X,\mathrm{soft}})-b_L)^2}{c_L}\right)} -\exp{\left(\frac{(\ln(\log(L_{X,\mathrm{soft, mean}})-b_L)^2}{c_L}\right)}\right]} \,,
    \end{split}
\end{equation}
with $a_L = -1.947\cdot 10^{17}$, $b_L = -1.572\cdot 10^{-4}$, $c_L = -0.2866$, the soft component of the X-ray luminosity is
\begin{equation}\label{eq:lxsoft}
    L_{X,\mathrm{soft}} = 10^{0.95 \log{(L_X)}+1.19} \,,
\end{equation}
and $L_{X,\mathrm{soft, mean}}$ is the soft component for a star with total X-ray luminosity given by the observational relation between stellar mass and mean X-ray luminosity \citep{2007A&A...468..353G}
\begin{equation} \label{eq:LxMstar}
    \log{L_{X, \mathrm{mean}}} = 1.54 \log(M_\star) + 30.31
\end{equation}
The consistency of the surface mass loss rate distribution however is reassuring, as it demonstrates that existing mass-loss prescriptions remain valid for compact discs, provided they are appropriately scaled to the total wind rate according to the relation above. In the following section, we normalize the mass-loss rates to the value of the cumulative mass-loss rate at the critical radius.
\begin{figure}
	\includegraphics{figures/Fig2.png}
	\caption{Cumulative mass loss rate as a function of radius for different critical radii.}
        \script{fig2.py}
	\label{fig:mass_loss_rate_cumulative}
\end{figure}

\subsection{Single disc evolution}

We investigated the disc long-term evolution using our 1-D viscous evolution code \textsc{spock} for discs with $r_1 = 5, 20,$ and $80$ au, $M_\mathrm{d} = 0.01\, M_\odot$, $\alpha=10^{-2}$ around a \SI{1}{M_\odot} star with $L_X = 2.29\cdot10^{30}$ erg/s, evolving the surface density according to equations~\ref{eq:visc}-\ref{eq:surf2}, \ref{eq:mdotnorm}-\ref{eq:LxMstar}.
In Figure~\ref{fig:photoevaporation_scenarios}, the left-hand panels show the surface density evolution obtained using our new, radius-scaled mass-loss prescription, while the right-hand panels display the corresponding results using the standard, unscaled prescription.
\begin{figure}
    \includegraphics[width=8cm]{figures/Fig3.png}
    \caption{Surface density evolution for discs with a critical radius of 5 au (top row), 20 au (middle row) and 80 au (bottom row) for the old prescription (right column) and the critical radius dependance (left column). The line colours indicate the different evolutionary stage as explained in the legend on top.}
    \script{fig3.py}
    \label{fig:photoevaporation_scenarios}
\end{figure}
The two approaches lead to strikingly different evolutionary outcomes. With the unscaled prescription, compact discs tend to disperse from the outside-in, gradually eroding their outer edges until the disc vanishes. In contrast, when the radius dependence is included, the discs evolve via inside-out dispersal, with gap opening driving the final dispersal. The difference in global lifetimes is also substantial, and the scaled prescription predicts significantly longer dispersal for the most compact discs.
This distinction is not only physically motivated, but also observationally significant. The inside-out dispersal predicted by the scaled prescription is consistent with the colour–colour evolution of young stellar objects (YSOs) observed in large surveys \citep{2013MNRAS.428.3327K}, which strongly disfavour outside-in clearing. The improved treatment of compact discs therefore provides a natural explanation for the observational evidence, reconciling theory and data.

\subsection{Population synthesis models}

We performed a population synthesis of $N_{\rm disc} = 10,000$ discs, with initial critical radii and masses drawn from the observationally derived distribution described in Section~\ref{sec:methods}. The resulting population naturally reproduces the observed disc fraction as a function of cluster age \citep[e.g.][]{2009AIPC.1158....3M} without the need of tuning the parameters as shown in Figure~\ref{fig:disc_fraction_age}. We scaled the initial disc fraction to be 86\%, considering that a fraction of discs would not form due to close binary interaction \citep[see e.g.][]{2001ApJ...553L.153H}. The resulting e-folding time for the critical radius-dependent population synthesis is $2.99 \pm 0.12$ Myr, which is compatible with the observationally derived one ($2.96 \pm 0.43$ Myr) and considerably larger than the previous prescription ($1.64 \pm 0.07$ Myr).
We then compared the accretion rate distributions as a function of age predicted by our population synthesis models with observational data from several stellar clusters, including Taurus, Lupus, Upper Scorpius, Chameleon I, Chameleon II, $\rho$ Ophiuchi, Corona Australis, Orionis Cloud A, Orionis Cloud B, $\gamma$ Vel, $\sigma$ Orionis, OriOB1a, OriOB1b \citep{2023ASPC..534..539M, 2022yCat..51630074T}. We evaluated the predictions obtained using both the standard (unscaled) prescription and the new, radius-scaled formulation of the photoevaporation-driven mass-loss rate.
The two approaches yield systematically different evolutionary trends. As shown in Fig.~\ref{fig:mass_accretion_age}, the unscaled prescription (panel a) results in a more rapid decline of accretion activity, underpredicting the number of systems with high accretion rates at ages beyond $\sim 4$ Myr. In contrast, the new radius-dependent prescription (panel b) produces a broader distribution of accretion rates and a more extended tail of strongly accreting discs persisting to older ages. This tail is a natural outcome of the weaker dependence of the dispersal timescale on stellar and disc parameters when the critical radius is accounted for explicitly.
The resulting accretion-rate–age distribution from the scaled model aligns more closely with the observations, which consistently reveal a minority population of long-lived, highly accreting discs in older clusters. Furthermore, the peak in the accretion rate is now shifted from $10^{-9}$ to $10^{-12}$, which is below our current detection threshold \citep{2023ApJ...944...90T} but it could be an interesting observable for future campaigns. This improved agreement suggests that incorporating the radius dependence into the photoevaporation prescription captures an essential physical ingredient missing from the standard approach, namely, the coupling between disc structure and the location of the flow launching region, which regulates both the efficiency and timescale of disc dispersal.

\begin{figure}
    \begin{centering}
	\includegraphics{figures/Fig4.png}
	\caption{Disc fraction as a function of time for the old (gray solid line) and new critical radius prescription (black solid line). The observational data from \citet{2009AIPC.1158....3M} and their exponential fit are shown with red dots and a red dashed line respectively. The e-folding time is overplotted with vertical dotted lines for each model and the observations.}
        \script{fig4.py}
	\label{fig:disc_fraction_age}
    \end{centering}
\end{figure}
\begin{figure*}
    \begin{centering}
        \includegraphics{figures/Fig5.png}
        \caption{Mass accretion rate as a function of time for the old prescription (panel a) and the critical radius criteria (panel b). The corner plots show the density distribution in the accretion rates and disc age, and the difference as a gray shadow region. The disc observed properties \citep{2023ASPC..534..539M, 2022yCat..51630074T} are overplotted as box plots.}
        \label{fig:mass_accretion_age}
        \script{fig5.py}
    \end{centering}
\end{figure*}
%\begin{figure*}
%	\centering
%	\includegraphics[width=17cm]{figures/Fig6_boxplot_kde.png}
%	\caption{Mass accretion rate as a function of time for the old prescription (panel a) and the critical radius criteria (panel b). The corner plots show the density distribution in the accretion rates and disc age, and the difference as a gray shadow region. The disc observed properties \citep{2023ASPC..534..539M, 2022yCat..51630074T} are overplotted as box plots.}
%	\label{fig:mass_accretion_age}
%\end{figure*}

\section{Conclusions}\label{sec:conclusions}

We have shown that compact discs, with a critical radius comparable with the gravitational radius for photoevaporation, follow the same local mass-loss profile as extended discs, but their integrated wind rates are reduced according to the previously found cumulative mass-loss rate distribution, given in equations~\ref{eq:mdotnorm}, \ref{eq:mdot},  \ref{eq:mdotmstar}, \ref{eq:mdotlx}, \ref{eq:lxsoft}.
Incorporating this scaling into viscous evolution models leads to longer disc lifetimes and inside-out dispersal, in agreement with observational constraints from colour–colour diagrams of young stellar populations. Population synthesis demonstrates that the corrected prescription reproduces both observed disc fractions and the presence of highly accreting discs at late ages ($t > 4$ Myr). The dispersal of compact discs cannot be captured by standard unscaled models, and radius-dependent photoevaporation prescriptions are essential to connect disc evolution theory with observations.

\section*{Acknowledgements}

We acknowledge the support of the DFG (German Research Foundation), Research Unit "Transition discs" - 325594231 and of the Excellence Cluster ORIGINS - EXC-2094 - 390783311. We acknowledge support from the Max Planck Society.

%%%%%%%%%%%%%%%%%%%%%%%%%%%%%%%%%%%%%%%%%%%%%%%%%%
\section*{Data Availability}

The data underlying this article will be shared on reasonable request to the corresponding author.

%%%%%%%%%%%%%%%%%%%% REFERENCES %%%%%%%%%%%%%%%%%%

\bibliographystyle{mnras}
\bibliography{biblio} % if your bibtex file is called example.bib

\appendix

% Don't change these lines
\bsp	% typesetting comment
\label{lastpage}
\end{document}

% End of mnras_template.tex
